\chapter{Complex curves and singular calibrated minimizers}\label{sec:calibrated-complex-curves}
Complex curves supply explicit mass-minimizing currents with singularities in higher codimension. We prove the complex-dimension-one calibration directly, following \cite{ChMin}, Example 16.5, p.~46, and develop the cusp mentioned in \cite{Minter}, p.~20. No higher-dimensional Wirtinger inequality or general classification of singular complex varieties is used.

\begin{lemma}[The Euclidean complex two-form has comass one]\label{lem:complex-two-form}
On $\mathbb C^k=\mathbb R^{2k}$ with its standard complex structure $J$, let
\[
 \omega=\sum_{j=1}^k dx_j\wedge dy_j.
\]
Then $d\omega=0$, $\|\omega\|_*=1$, and for an orthonormal pair $(v,w)$,
\[
 \omega(v,w)=1\quad\Longleftrightarrow\quad w=Jv.
\]
Thus equality selects complex lines with their complex orientation.
\end{lemma}
\begin{proof}
The coefficients are constant, so the form is closed. Direct calculation gives $\omega(v,w)=\langle Jv,w\rangle$. The Cauchy--Schwarz inequality bounds its absolute value by one; the value one occurs at $(v,Jv)$, since $J$ is orthogonal and $v\perp Jv$. Equality with positive sign in Cauchy--Schwarz is exactly $w=Jv$. This is the two-dimensional calculation in \cite{ChMin}, Example 16.5.
\end{proof}

\begin{proposition}[Area and orientation of a holomorphic parametrization]\label{prop:holomorphic-area-identity}
Let $F:D\subset\mathbb C\to\mathbb C^k$ be holomorphic, and write the domain coordinate as $z=s+it$. Then
\[
 F_t=JF_s,\qquad J_2DF=|F'(z)|^2,
 \qquad F^*\omega=|F'(z)|^2\,ds\wedge dt.
\]
At every point with $F'\ne0$, the oriented tangent plane is a positively oriented complex line.
\end{proposition}
\begin{proof}
The Cauchy--Riemann equations give the first identity componentwise. The two derivative vectors are orthogonal and have equal length $|F'|$, so their Gram determinant is $|F'|^4$. The area Jacobian is its square root. Lemma \ref{lem:complex-two-form} gives $\omega(F_s,F_t)=|F_s|^2$. At a critical point both sides vanish, so the formulas remain valid there.
\end{proof}

\begin{proposition}[Holomorphic disks give calibrated integral currents]\label{prop:holomorphic-current-minimizer}
Let $D\subset\mathbb C$ be a bounded domain with smooth boundary, and suppose $F$ is holomorphic on a neighborhood of $\overline D$. Orient $D$ by $ds\wedge dt$ and put $T=F_\#\llbracket D\rrbracket$. Then $T$ is a compactly supported integral current,
\[
 \partial T=F_\#\llbracket\partial D\rrbracket,
 \qquad \mathbf M(T)=T(\omega)=\int_D|F'|^2\,ds\,dt.
\]
It minimizes mass among compactly supported finite-mass real two-currents in $\mathbb C^k$ with the same boundary. The formula counts multiplicity if $F$ is not injective.
\end{proposition}
\begin{proof}
The map is smooth and Lipschitz on a neighborhood of the compact closure; Theorem \ref{thm:lipschitz-current-pushforward} gives integrality and the boundary formula. The pushforward representation gives $\mathbf M(T)\leq\int_DJ_2DF$. Proposition \ref{prop:holomorphic-area-identity} gives $T(\omega)=\int_DJ_2DF$, whereas the comass bound gives $T(\omega)\leq\mathbf M(T)$. Equality follows, showing in particular that orientation cancellation cannot reduce the mass. Corollary \ref{cor:calibrated-euclidean-boundary} gives the minimizing property. Compactly supported cutoff forms justify each evaluation.
\end{proof}

\begin{example}[The complex cusp as an integral minimizer]\label{ex:complex-cusp-current}
For $D=\{z\in\mathbb C:|z|<1\}$ and $F(z)=(z^2,z^3)$, the current $T=F_\#\llbracket D\rrbracket$ is calibrated and mass minimizing with its prescribed outer boundary. The map is injective, is an immersion away from zero, and has area Jacobian
\[
 J_2DF(z)=4|z|^2+9|z|^4.
\]
Its image satisfies $w^2=v^3$ in coordinates $(v,w)\in\mathbb C^2$. The origin is not asserted to be a smooth complex submanifold point.
\end{example}
\begin{proof}
If two nonzero parameters have the same image, their ratio has both square and cube equal to one, so the parameters agree. Zero has only the zero preimage. The derivative $(2z,3z^2)$ vanishes exactly at zero. Its squared norm gives the displayed Jacobian by Proposition \ref{prop:holomorphic-area-identity}; Proposition \ref{prop:holomorphic-current-minimizer} supplies integrality and minimization. The boundary is the smooth embedded outer curve $F(\partial D)$, since the derivative is nonzero there. The algebraic equation follows by substitution. Compare the cusp example of \cite{Minter}, p.~20; the parameter convention here interchanges the source's two coordinates.
\end{proof}

\begin{proposition}[Exact density of the calibrated cusp]\label{prop:cusp-density}
Let $V$ be the rectifiable varifold associated to the cusp current in Example \ref{ex:complex-cusp-current}. For $0<r<\sqrt2$, let $R\in(0,1)$ solve $R^4+R^6=r^2$. Then
\[
 \|V\|(B_r(0))=2\pi R^4+3\pi R^6,
 \qquad\frac{\|V\|(B_r(0))}{\pi r^2}
 =\frac{2+3R^2}{1+R^2}.
\]
In particular $\Theta^2(\|V\|,0)=2$.
\end{proposition}
\begin{proof}
The norm squared of $F(z)$ is $|z|^4+|z|^6$, strictly increasing in $|z|$. Thus the preimage of the ambient ball is the parameter disk of radius $R$. Injectivity and the area formula give
\[
 \|V\|(B_r(0))=2\pi\int_0^R(4t^2+9t^4)t\,dt
 =2\pi R^4+3\pi R^6.
\]
Divide by $\pi(R^4+R^6)$ and let $R\downarrow0$. The critical point contributes no area atom.
\end{proof}

\begin{proposition}[The cusp has a unique double-plane tangent]\label{prop:cusp-tangent}
For the cusp varifold $V$, the full family of dilations $V_{0,r}$ converges as $r\downarrow0$ to $2|P|$, where $P=\mathbb C\times\{0\}\subset\mathbb C^2$.
\end{proposition}
\begin{proof}
In the dilation integral make the substitution $z=\sqrt r\,\zeta$. The position is then $(\zeta^2,\sqrt r\,\zeta^3)$, the plane away from zero is the complex span of $(2\zeta,3\sqrt r\,\zeta^2)$, and the rescaled area factor is
\[
 4|\zeta|^2+9r|\zeta|^4.
\]
For a continuous Grassmann-bundle test supported over $B_A(0)$, the contributing parameters satisfy $|\zeta|\leq\sqrt A$, because the first coordinate is $\zeta^2$. On this fixed bounded disk the area factors are uniformly bounded for $r\leq1$. Except at the zero parameter the planes converge to $P$, so dominated convergence gives the limit integral with position $(\zeta^2,0)$ and factor $4|\zeta|^2$. The planar area formula for $\zeta\mapsto\zeta^2$ identifies this with twice the integral over $P$; every nonzero image point has two preimages. The calculation holds along the full limit $r\downarrow0$, proving uniqueness as a varifold.
\end{proof}

\begin{corollary}[A singular calibrated surface with a flat tangent]\label{cor:cusp-flat-singular}
The cusp varifold is stationary away from its outer boundary, smooth at every nonzero interior point, and singular at zero. Its origin is a flat singular point of density two with a unique tangent cone. Thus higher-codimension mass minimization permits this type of flat singularity.
\end{corollary}
\begin{proof}
Stationarity follows from Proposition \ref{prop:calibrated-current-stationary}. Injectivity and immersion give smoothness locally away from zero. If the support were a smooth embedded surface near zero, Proposition \ref{prop:cusp-tangent} would make its tangent plane horizontal, so it would be a single graph over $P$ on a smaller neighborhood. For each sufficiently small nonzero $v$, however, the two parameters solving $z^2=v$ give opposite nonzero values of $w=z^3$. Both points lie in that neighborhood, contradicting the graph property. Propositions \ref{prop:cusp-density} and \ref{prop:cusp-tangent} give the stated density and tangent. This is a concrete example under Definition \ref{def:flat-singular-point}, not a classification of all flat singularities.
\end{proof}
